\section{Quá trình thiết kế}

\subsection{Kiến trúc thông tin}
Kiến trúc thông tin (Information Architecture -- IA) của hệ thống được tổ chức tinh gọn theo mô hình phân cấp phẳng, chuẩn Mobile-first nhằm giảm thiểu tối đa số lần nhấp chuột (clicks) và độ sâu điều hướng. Ứng dụng bao gồm 4 phân hệ chính gắn liền với thanh điều hướng đáy (Bottom Navigation Bar) cố định:
\begin{enumerate}
    \item \textbf{Trang chủ (Dashboard / Home)}: Tổng quan lịch hẹn sắp tới, lối tắt đặt lịch nhanh một chạm (Quick Rebook), thông báo mới nhất và tình trạng nhanh của các thú cưng.
    \item \textbf{Đặt lịch (Booking Matrix)}: Luồng chọn thú cưng, danh mục dịch vụ spa/cắt tỉa, ma trận khung giờ khả dụng theo ngày và màn hình xác nhận chi tiết đính kèm ghi chú y tế.
    \item \textbf{Tiến độ (Live Tracking)}: Trung tâm theo dõi tiến độ thời gian thực 4 mốc (\textit{Đã nhận} $\rightarrow$ \textit{Đang chăm sóc} $\rightarrow$ \textit{Hoàn tất} $\rightarrow$ \textit{Chờ đón}), tích hợp nút gửi thêm dặn dò bổ sung và mã QR nhận thú cưng khi hoàn thành.
    \item \textbf{Hồ sơ (Pet Profiles \& History)}: Quản lý danh sách thú cưng, hồ sơ y tế chuyên sâu (tiền sử dị ứng, sổ tiêm chủng), lịch sử dịch vụ cá nhân hóa, hóa đơn điện tử và biểu đồ theo dõi chi tiêu hàng tháng.
\end{enumerate}

\subsection{Luồng thao tác người dùng}
Dựa trên việc tái cấu trúc quy trình cũ, nhóm đã thiết kế Quy trình đề xuất (Quy trình To-Be) gồm 6 bước khép kín, tự động hóa và minh bạch thông tin:
\begin{enumerate}
    \item \textbf{Bước 1 -- Đăng nhập và chọn thông tin}: Chủ nuôi mở ứng dụng, chọn thú cưng, gói dịch vụ mong muốn và trực tiếp lựa chọn khung giờ trống trên ma trận thời gian thực (các khung giờ đã kín tự động bị khóa).
    \item \textbf{Bước 2 -- Xác nhận tức thì \& Đính kèm cảnh báo}: Hệ thống ghi nhận lịch hẹn ngay lập tức, tự động trích xuất các dặn dò dị ứng/thuốc từ hồ sơ thú cưng để đính kèm vào đơn, đồng thời sinh mã định danh và xác nhận trên ứng dụng.
    \item \textbf{Bước 3 -- Tiếp nhận tại tiệm qua mã QR}: Khi đưa thú cưng đến tiệm, chủ nuôi xuất trình mã QR tiếp nhận; nhân viên quét mã và màn hình tiếp nhận lập tức hiển thị cảnh báo dị ứng màu đỏ nổi bật để kỹ thuật viên nắm rõ.
    \item \textbf{Bước 4 -- Theo dõi tiến độ 4 mốc thời gian thực}: Trong suốt quá trình thú cưng được chăm sóc, hệ thống hiển thị mốc trạng thái trực quan và chủ động gửi thông báo đẩy đến điện thoại của chủ nuôi mỗi khi hoàn thành một công đoạn.
    \item \textbf{Bước 5 -- Nhận thông báo đón tức thời}: Ngay khi dịch vụ hoàn tất và chuyển sang mốc ''Chờ đón'', chủ nuôi nhận được thông báo kèm hình ảnh đối chiếu trước/sau spa để chủ động di chuyển đến tiệm đón bé.
    \item \textbf{Bước 6 -- Lưu trữ lịch sử chăm sóc tự động}: Toàn bộ thông tin về dịch vụ, loại sữa tắm đã dùng, hình ảnh kết quả và hóa đơn điện tử được tự động lưu vào hồ sơ thú cưng, sẵn sàng cho việc tra cứu hoặc tái đặt lịch (Rebook) ở các lần sau.
\end{enumerate}

\subsection{Phác thảo và phát triển ý tưởng}
Quá trình phát triển ý tưởng bắt đầu từ các bản phác thảo nhanh trên giấy (Sketches \& Crazy Eights) nhằm khám phá nhiều cách bố trí thành phần giao diện khác nhau:
\begin{itemize}
    \item Thử nghiệm cách hiển thị lịch hẹn: từ dạng lịch tháng truyền thống (gây khó nhìn trên màn hình hẹp) chuyển sang dạng ma trận thẻ thời gian theo ngày (Timeslot Matrix Chips) giúp chọn giờ bằng một chạm.
    \item Thử nghiệm hiển thị cảnh báo y tế: từ dạng hộp thoại cảnh báo (modal popup) dễ bị người dùng lướt qua sang dạng thẻ cảnh báo tương phản cao gắn cố định vào tiêu đề đơn tiếp nhận.
    \item Thử nghiệm thanh tiến độ: từ dạng thanh phần trăm trừu tượng (0--100\%) chuyển sang mô hình 4 mốc trạng thái cụ thể gắn liền với hành động thực tế tại cơ sở spa.
\end{itemize}

\subsection{Bảng phân cảnh}
Nhóm đã trực quan hóa câu chuyện trải nghiệm của người dùng qua Bảng phân cảnh (Storyboard) 6 khung tranh chuẩn Expressive Stick-figure UI, thể hiện sự chuyển biến cảm xúc từ lo âu, bất an trong quy trình cũ sang an tâm, chủ động và hài lòng trong hệ thống mới:

\begin{figure}[H]
\centering
\includegraphics[width=0.90\textwidth, keepaspectratio]{storyboard.png}
\caption{Bảng phân cảnh trải nghiệm -- Kịch bản đặt lịch và theo dõi tiến độ của Nguyễn Hoàng Lan}
\label{fig:storyboard_1}
\end{figure}

\subsection{Khung giao diện}
Nhóm đã thiết kế bộ Khung giao diện toàn diện gồm 32 màn hình Wireframe chuẩn Mobile-first ($430 \times 932\text{px}$) tích hợp đầy đủ 5 trạng thái giao diện chuẩn công nghiệp:
\begin{itemize}
    \item \textbf{Cụm màn hình Trang chủ \& Đặt lịch (01--07)}: Bảng điều khiển trung tâm, cài đặt thông báo, danh mục dịch vụ, ma trận chọn giờ trống, đặt lịch nhiều thú cưng song song và màn hình xác nhận thành công.
    \item \textbf{Cụm màn hình Hồ sơ \& Y tế (08--10)}: Quản lý hồ sơ thú cưng, hồ sơ cảnh báo dị ứng y tế và sổ tiêm chủng điện tử.
    \item \textbf{Cụm màn hình Tiếp nhận \& Bàn giao 2 lớp (11--13)}: Mã QR check-in tiếp tân, biên bản bàn giao tình trạng thú cưng và xác nhận an toàn 2 lớp.
    \item \textbf{Cụm màn hình Theo dõi tiến độ thời gian thực (14--19)}: 4 màn hình chi tiết cho 4 mốc tiến độ (Đã nhận, Đang chăm sóc, Hoàn tất, Chờ đón), theo dõi song song nhiều bé và hình ảnh phòng cách ly an toàn.
    \item \textbf{Cụm màn hình Trả thú cưng \& Lịch sử dịch vụ (20--27)}: Mã QR đối soát trả bé, báo cáo hình ảnh trước/sau spa, đánh giá 5 sao, trung tâm lịch sử dịch vụ, chi tiết phiên chăm sóc, đặt lại nhanh (Rebook), theo dõi ngân sách và hóa đơn điện tử.
    \item \textbf{Cụm màn hình 5 trạng thái giao diện}: Trạng thái mặc định (Default), Trạng thái rỗng (Empty), Trạng thái đang tải (Loading), Trạng thái lỗi mạng/xung đột dị ứng (Error \& Conflict), và Trạng thái hoàn tất thành công (Success).
\end{itemize}

\subsection{Các phương án thiết kế}
Trong quá trình thực hiện đồ án, nhóm áp dụng nguyên tắc thiết kế tập trung chuyên sâu (Focused Single-Architecture Design Approach):
\begin{itemize}
    \item \textbf{Phương án thiết kế tối ưu duy nhất}: Thay vì phân tán nguồn lực vào các phương án thử nghiệm kém tối ưu hoặc tạo ra các kịch bản so sánh hình thức (A, B, C), nhóm tập trung toàn bộ nguồn lực để hoàn thiện chuyên sâu \textbf{một phương án thiết kế tối ưu duy nhất} được xây dựng trực tiếp từ kết quả nghiên cứu người dùng và Proposal.
    \item \textbf{Lý giải quyết định thiết kế}: Phương án tối ưu này giải quyết trực diện và triệt để 4 điểm đau cốt lõi:
    \begin{enumerate}
        \item Sử dụng ma trận chọn giờ trống thời gian thực thay vì lịch thả rơi (Dropdown) giúp giảm $65\%$ thao tác nhầm lẫn.
        \item Cơ chế tự động đính kèm thông tin y tế từ hồ sơ loại bỏ $100\%$ nguy cơ bỏ sót dặn dò nguy hiểm.
        \item Thanh tiến trình 4 mốc chuẩn hóa kết hợp thông báo đẩy giải quyết trọn vẹn sự bất an của chủ nuôi trong giờ làm việc.
        \item Khung lưu trữ lịch sử cá nhân hóa tích hợp tính năng đặt lại nhanh (Rebook) trong 1 chạm mang lại sự tiện lợi tối đa cho các lần chăm sóc định kỳ.
    \end{enumerate}
\end{itemize}
